\section{Основная часть}

\subsection{Создание представления}

В базе данных одно логическое начертание может быть представлено несколькими
строками таблицы \mintinline{text}{Typeface}: отдельные строки описывают его
варианты формата, насыщенности и наклона. Логическое начертание однозначно
задаётся идентификатором гарнитуры \mintinline{sql}{id_font_family} и названием
\mintinline{sql}{name}.

Для каждого начертания требуется определить два показателя:

\begin{itemize}
    \item \mintinline{sql}{format_count}~--- число различных поддерживаемых
    форматов;
    \item \mintinline{sql}{weight_count}~--- число различных градаций
    насыщенности.
\end{itemize}

Операция \mintinline{sql}{COUNT(DISTINCT ...)} исключает повторный учёт одного
формата или одной насыщенности. Поле \mintinline{sql}{id_font_family}
сохраняется в представлении для последующего соединения начертаний с
гарнитурами и категориями. Команда создания представления приведена в
листинге~\ref{lst:create-view}.

\codeListing[sql]{../src/create.sql}{Создание представления
\mintinline{sql}{typeface_format_weight_stats}}{\label{lst:create-view}}

В терминах реляционной алгебры построение представления можно записать в виде

\[
V = \gamma_{\substack{id\_font\_family,\ name;\\
\operatorname{count\_distinct}(id\_format) \rightarrow format\_count,\\
\operatorname{count\_distinct}(id\_weight) \rightarrow weight\_count}}
(Typeface).
\]

Использование \mintinline{sql}{CREATE OR REPLACE VIEW} позволяет повторно
выполнить скрипт после изменения определения представления. Обычное
представление не хранит отдельную копию строк: PostgreSQL вычисляет его
содержимое по базовой таблице при выполнении запроса.

\subsection{Проверка представления}

После выполнения команды создания представление появилось в схеме
\mintinline{text}{public}. Для проверки его содержимого был выполнен запрос из
листинга~\ref{lst:query-01}.

\codeListing[sql]{../src/query-01.sql}{Проверка содержимого представления}{\label{lst:query-01}}

Фрагмент полученного результата приведён в таблице~\ref{table:view-result}.
Для наглядности выбраны строки с наибольшим числом форматов и градаций
насыщенности.

\begin{table}[H]
    \centering
    \caption{Фрагмент содержимого представления}
    \label{table:view-result}
    \begin{tblr}{
        colspec={|r|l|r|r|},
        row{1}={font=\bfseries},
        hlines,
        vlines,
    }
        \mintinline{sql}{id_font_family} & \mintinline{sql}{typeface_name} &
        \mintinline{sql}{format_count} & \mintinline{sql}{weight_count} \\
        159 & Typeface 159 & 4 & 9 \\
        166 & Typeface 166 & 4 & 9 \\
        156 & Typeface 156 & 4 & 8 \\
        164 & Typeface 164 & 4 & 8 \\
        167 & Typeface 167 & 4 & 8 \\
        147 & Typeface 147 & 4 & 7 \\
    \end{tblr}
\end{table}

Результат подтверждает, что представление содержит по одной строке для
логического начертания и обе требуемые агрегатные характеристики. В
сгенерированном наборе данных число форматов изменяется от 1 до 4, а число
градаций насыщенности~--- от 2 до 9.

\clearpage
\subsection{Запрос с условием}

Пусть $A=2$. Требуется найти количество начертаний, поддерживающих более двух
форматов. Для этого к строкам представления применяется условие
\mintinline{sql}{format_count > 2}, после чего оставшиеся строки подсчитываются.
Запрос приведён в листинге~\ref{lst:query-02}.

\codeListing[sql]{../src/query-02.sql}{Подсчёт начертаний с числом форматов больше
$A$}{\label{lst:query-02}}

Результат выполнения запроса:

\begin{table}[H]
    \centering
    \caption{Количество начертаний, поддерживающих более двух форматов}
    \label{table:query-02-result}
    \begin{tblr}{colspec={|c|}, row{1}={font=\bfseries}, hlines, vlines}
        \mintinline{sql}{typeface_count} \\
        39 \\
    \end{tblr}
\end{table}

Таким образом, условию соответствует 39 логических начертаний.

\subsection{Охват форматов по категориям}

Для определения категории каждого начертания представление соединяется с
таблицей гарнитур \mintinline{text}{Font_family}, а затем с таблицей категорий
\mintinline{text}{Category}. После соединения значения
\mintinline{sql}{format_count} суммируются по категории. Запрос приведён в
листинге~\ref{lst:query-03}.

\codeListing[sql]{../src/query-03.sql}{Первые пять категорий по охвату
поддерживаемых форматов}{\label{lst:query-03}}

В данном запросе представление используется как источник уже рассчитанного
числа форматов каждого начертания. Это исключает необходимость повторять
группировку таблицы \mintinline{text}{Typeface} в тексте итогового запроса.
Результат приведён в таблице~\ref{table:query-03-result}.

\begin{table}[H]
    \centering
    \caption{Категории с наибольшим охватом форматов}
    \label{table:query-03-result}
    \begin{tblr}{
        colspec={|r|X[l]|r|},
        row{1}={font=\bfseries},
        hlines,
        vlines,
    }
        \mintinline{sql}{id_category} & Категория & Охват форматов \\
        10 & Neo-grotesque & 56 \\
        9 & Grotesque & 54 \\
        2 & Sans-serif & 38 \\
        11 & Transitional & 33 \\
        3 & Monospace & 23 \\
    \end{tblr}
\end{table}

Наибольший суммарный охват имеет категория
\mintinline{text}{Neo-grotesque}: входящие в неё начертания в совокупности
поддерживают 56 вариантов формата.

\subsection{План выполнения итогового запроса}

План был получен командой \mintinline{sql}{EXPLAIN (ANALYZE, BUFFERS, FORMAT JSON)}.
PostgreSQL разворачивает определение представления, последовательно читает
таблицу начертаний и группирует её по гарнитуре и названию. Полученный набор
соединяется с таблицей гарнитур с помощью хеш-соединения. Категории находятся
по первичному ключу; узел \mintinline{text}{Memoize} повторно использует уже
найденные строки категорий. После этого выполняется агрегирование по категории
и сортировка \mintinline{text}{top-N heapsort}, достаточная для получения пяти
первых строк.

Для актуального набора данных время планирования составило 0,414~мс, время
выполнения~--- 1,489~мс. JSON-планы всех запросов сохранены отдельно в каталоге
\mintinline{text}{practice-01/explain} и могут быть загружены в приложение для
визуализации планов PostgreSQL.
